% --- Archon physics formalization source begin ---
% archon:physics
% archon:covers PhyXMiniProblems/problem_phyx_mini_0778.lean
% archon:source-report reports/phyx_mini/problem_phyx_mini_0778.source.json

\chapter{Physics problem phyx\_mini\_0778}
\label{ch:PhyXMiniProblems_problem_phyx_mini_0778}

\paragraph{Problem source.}
\#\# Physical scenario

A \textbackslash{}(12.0\textbackslash{}text\{ kg\}\textbackslash{}) box resting on a horizontal, frictionless surface is attached to a \textbackslash{}(5.00\textbackslash{}text\{ kg\}\textbackslash{}) weight by a thin, light wire that passes over a frictionless pulley as shown in figure. The pulley has the shape of a uniform solid disk of mass \textbackslash{}(2.00\textbackslash{}text\{ kg\}\textbackslash{}) and diameter \textbackslash{}(0.500\textbackslash{}text\{ m\}\textbackslash{}). After the system is released.

\#\# Question

Find the acceleration of the box

\#\# Answer choices

- A: A: \textbackslash{}( 1.36 \textbackslash{}text\{ m/s\}\textasciicircum{}2 \textbackslash{})

- B: B: \textbackslash{}( 2.72 \textbackslash{}text\{ m/s\}\textasciicircum{}2 \textbackslash{})

- C: C: \textbackslash{}( 1.72 \textbackslash{}text\{ m/s\}\textasciicircum{}2 \textbackslash{})

- D: D: \textbackslash{}( 4.97 \textbackslash{}text\{ m/s\}\textasciicircum{}2 \textbackslash{})

\#\# Figure caption (auxiliary; use the image as primary evidence)

The image shows a physics setup involving a pulley system:

1. **Objects**:
   - **Two blocks**: One block labeled "12.0 kg" on the left, which is resting on a flat horizontal surface. The second block, labeled "5.00 kg", is hanging off the edge of the surface.
   - **Pulley**: A fixed pulley is attached to the edge of the horizontal surface via a clamp.
  
2. **Connections**:
   - A rope or cable is connected to both blocks and passes over the pulley. The 12.0 kg block pulls horizontally, and the 5.00 kg block hangs vertically.

3. **Text**:
   - The labels "12.0 kg" and "5.00 kg" indicate the mass of each block.

4. **Scene**:
   - The apparatus is set up with the heavier block on a horizontal table or surface, while the lighter block hangs off, connected and balanced by the pulley system. The clamp holds the pulley securely to the edge of the table.

This setup is commonly used to demonstrate principles of physics related to tension, gravity, and mechanical advantage.

\#\# Dataset metadata

Dynamics; Physical Model Grounding Reasoning, Multi-Formula Reasoning

\paragraph{Recorded answer/context.}
B: B: \textbackslash{}( 2.72 \textbackslash{}text\{ m/s\}\textasciicircum{}2 \textbackslash{})

\paragraph{Figure/image path.}
/root/proposal\_for\_physic/science-mango/phyx\_mini\_run/phyx\_data/test\_image/778.png

\paragraph{Formalization target.}
create a compiling Lean file with sorry bodies at `PhyXMiniProblems/problem\_phyx\_mini\_0778.lean`. The Lean declarations must preserve the physical quantities, dimensions or dimensional roles, figure labels, governing-law hypotheses, and final relation expressed by this problem.
Use Mathlib/Physlib names found through LeanExplore where available. If a physics API is missing, introduce faithful local abstractions rather than scalar placeholder aliases.

\begin{theorem}[Physics formalization target]
\label{thm:physics:phyx_mini_0778:target}
\lean{PhyXMiniProblems.ProblemPhyXMini0778.problem_phyx_mini_0778}
\uses{def:physics:phyx-mini-0778:phyxminiproblems-problemphyxmini0778-accelerationinmeterspersecondsquared, def:physics:phyx-mini-0778:phyxminiproblems-problemphyxmini0778-tabletoppulleysetup, def:physics:phyx-mini-0778:phyxminiproblems-problemphyxmini0778-matcheswrittenscenario, def:physics:phyx-mini-0778:phyxminiproblems-problemphyxmini0778-matchesprimaryfigure, def:physics:phyx-mini-0778:phyxminiproblems-problemphyxmini0778-matchesproblemreadouts, def:physics:phyx-mini-0778:phyxminiproblems-problemphyxmini0778-usesstandardnearearthgravity, def:physics:phyx-mini-0778:phyxminiproblems-problemphyxmini0778-hasphysicalparameters, def:physics:phyx-mini-0778:phyxminiproblems-problemphyxmini0778-satisfiesmassivepulleydynamics, def:physics:phyx-mini-0778:phyxminiproblems-problemphyxmini0778-recordedanswerchoice, def:physics:phyx-mini-0778:phyxminiproblems-problemphyxmini0778-matchesdisplayedanswer}
The assigned autoformalize agent should translate this physics problem into Lean declarations in the covered file, with theorem and lemma proof bodies written as `by sorry`.
\end{theorem}
\begin{proof}
This is an autoformalization task, not a proof task. Produce statements that can later be proved by the physics prover without weakening the physical model.
\end{proof}

% --- Archon declaration topology begin ---
\section{Declaration topology}
\begin{definition}[acceleration Dimension]
  \label{def:physics:phyx-mini-0778:phyxminiproblems-problemphyxmini0778-accelerationdimension}
  \lean{PhyXMiniProblems.ProblemPhyXMini0778.accelerationDimension}
  The physical dimension L T⁻² of a linear acceleration magnitude
\end{definition}
\begin{proof}
  This is the declared model definition of acceleration Dimension.
\end{proof}
\begin{definition}[force Dimension]
  \label{def:physics:phyx-mini-0778:phyxminiproblems-problemphyxmini0778-forcedimension}
  \lean{PhyXMiniProblems.ProblemPhyXMini0778.forceDimension}
  The physical dimension M L T⁻² of a force magnitude
\end{definition}
\begin{proof}
  This is the declared model definition of force Dimension.
\end{proof}
\begin{definition}[angular Acceleration Dimension]
  \label{def:physics:phyx-mini-0778:phyxminiproblems-problemphyxmini0778-angularaccelerationdimension}
  \lean{PhyXMiniProblems.ProblemPhyXMini0778.angularAccelerationDimension}
  Angular acceleration has dimension T⁻²; radians are dimensionless
\end{definition}
\begin{proof}
  This is the declared model definition of angular Acceleration Dimension.
\end{proof}
\begin{definition}[moment Of Inertia Dimension]
  \label{def:physics:phyx-mini-0778:phyxminiproblems-problemphyxmini0778-momentofinertiadimension}
  \lean{PhyXMiniProblems.ProblemPhyXMini0778.momentOfInertiaDimension}
  A scalar axial moment of inertia has dimension M L²
\end{definition}
\begin{proof}
  This is the declared model definition of moment Of Inertia Dimension.
\end{proof}
\begin{definition}[torque Dimension]
  \label{def:physics:phyx-mini-0778:phyxminiproblems-problemphyxmini0778-torquedimension}
  \lean{PhyXMiniProblems.ProblemPhyXMini0778.torqueDimension}
  A torque magnitude has dimension M L² T⁻²
\end{definition}
\begin{proof}
  This is the declared model definition of torque Dimension.
\end{proof}
\begin{definition}[Mass Quantity]
  \label{def:physics:phyx-mini-0778:phyxminiproblems-problemphyxmini0778-massquantity}
  \lean{PhyXMiniProblems.ProblemPhyXMini0778.MassQuantity}
  A nonnegative, unit-independent physical mass
\end{definition}
\begin{proof}
  This is the declared model definition of Mass Quantity.
\end{proof}
\begin{definition}[Length Quantity]
  \label{def:physics:phyx-mini-0778:phyxminiproblems-problemphyxmini0778-lengthquantity}
  \lean{PhyXMiniProblems.ProblemPhyXMini0778.LengthQuantity}
  A nonnegative, unit-independent physical length
\end{definition}
\begin{proof}
  This is the declared model definition of Length Quantity.
\end{proof}
\begin{definition}[Acceleration Quantity]
  \label{def:physics:phyx-mini-0778:phyxminiproblems-problemphyxmini0778-accelerationquantity}
  \lean{PhyXMiniProblems.ProblemPhyXMini0778.AccelerationQuantity}
  \uses{def:physics:phyx-mini-0778:phyxminiproblems-problemphyxmini0778-accelerationdimension}
  A nonnegative, unit-independent linear acceleration magnitude
\end{definition}
\begin{proof}
  This is the declared model definition of Acceleration Quantity.
\end{proof}
\begin{definition}[Force Quantity]
  \label{def:physics:phyx-mini-0778:phyxminiproblems-problemphyxmini0778-forcequantity}
  \lean{PhyXMiniProblems.ProblemPhyXMini0778.ForceQuantity}
  \uses{def:physics:phyx-mini-0778:phyxminiproblems-problemphyxmini0778-forcedimension}
  A nonnegative, unit-independent force magnitude
\end{definition}
\begin{proof}
  This is the declared model definition of Force Quantity.
\end{proof}
\begin{definition}[Angular Acceleration Quantity]
  \label{def:physics:phyx-mini-0778:phyxminiproblems-problemphyxmini0778-angularaccelerationquantity}
  \lean{PhyXMiniProblems.ProblemPhyXMini0778.AngularAccelerationQuantity}
  \uses{def:physics:phyx-mini-0778:phyxminiproblems-problemphyxmini0778-angularaccelerationdimension}
  A nonnegative angular-acceleration magnitude about the pulley axle
\end{definition}
\begin{proof}
  This is the declared model definition of Angular Acceleration Quantity.
\end{proof}
\begin{definition}[Moment Of Inertia Quantity]
  \label{def:physics:phyx-mini-0778:phyxminiproblems-problemphyxmini0778-momentofinertiaquantity}
  \lean{PhyXMiniProblems.ProblemPhyXMini0778.MomentOfInertiaQuantity}
  \uses{def:physics:phyx-mini-0778:phyxminiproblems-problemphyxmini0778-momentofinertiadimension}
  A nonnegative scalar moment of inertia about the pulley axle
\end{definition}
\begin{proof}
  This is the declared model definition of Moment Of Inertia Quantity.
\end{proof}
\begin{definition}[Torque Quantity]
  \label{def:physics:phyx-mini-0778:phyxminiproblems-problemphyxmini0778-torquequantity}
  \lean{PhyXMiniProblems.ProblemPhyXMini0778.TorqueQuantity}
  \uses{def:physics:phyx-mini-0778:phyxminiproblems-problemphyxmini0778-torquedimension}
  A nonnegative torque magnitude about the pulley axle
\end{definition}
\begin{proof}
  This is the declared model definition of Torque Quantity.
\end{proof}
\begin{definition}[coherent SIReadout]
  \label{def:physics:phyx-mini-0778:phyxminiproblems-problemphyxmini0778-coherentsireadout}
  \lean{PhyXMiniProblems.ProblemPhyXMini0778.coherentSIReadout}
  Read a nonnegative dimensionful quantity in coherent SI base units
\end{definition}
\begin{proof}
  This is the declared model definition of coherent SIReadout.
\end{proof}
\begin{definition}[mass In Kilograms]
  \label{def:physics:phyx-mini-0778:phyxminiproblems-problemphyxmini0778-massinkilograms}
  \lean{PhyXMiniProblems.ProblemPhyXMini0778.massInKilograms}
  \uses{def:physics:phyx-mini-0778:phyxminiproblems-problemphyxmini0778-massquantity, def:physics:phyx-mini-0778:phyxminiproblems-problemphyxmini0778-coherentsireadout}
  Kilogram readout of a physical mass
\end{definition}
\begin{proof}
  This is the declared model definition of mass In Kilograms.
\end{proof}
\begin{definition}[length In Meters]
  \label{def:physics:phyx-mini-0778:phyxminiproblems-problemphyxmini0778-lengthinmeters}
  \lean{PhyXMiniProblems.ProblemPhyXMini0778.lengthInMeters}
  \uses{def:physics:phyx-mini-0778:phyxminiproblems-problemphyxmini0778-lengthquantity, def:physics:phyx-mini-0778:phyxminiproblems-problemphyxmini0778-coherentsireadout}
  Metre readout of a physical length
\end{definition}
\begin{proof}
  This is the declared model definition of length In Meters.
\end{proof}
\begin{definition}[acceleration In Meters Per Second Squared]
  \label{def:physics:phyx-mini-0778:phyxminiproblems-problemphyxmini0778-accelerationinmeterspersecondsquared}
  \lean{PhyXMiniProblems.ProblemPhyXMini0778.accelerationInMetersPerSecondSquared}
  \uses{def:physics:phyx-mini-0778:phyxminiproblems-problemphyxmini0778-accelerationquantity, def:physics:phyx-mini-0778:phyxminiproblems-problemphyxmini0778-coherentsireadout}
  Metres-per-second-squared readout of a linear acceleration magnitude
\end{definition}
\begin{proof}
  This is the declared model definition of acceleration In Meters Per Second Squared.
\end{proof}
\begin{definition}[force In Newtons]
  \label{def:physics:phyx-mini-0778:phyxminiproblems-problemphyxmini0778-forceinnewtons}
  \lean{PhyXMiniProblems.ProblemPhyXMini0778.forceInNewtons}
  \uses{def:physics:phyx-mini-0778:phyxminiproblems-problemphyxmini0778-forcequantity, def:physics:phyx-mini-0778:phyxminiproblems-problemphyxmini0778-coherentsireadout}
  Newton readout of a physical force magnitude
\end{definition}
\begin{proof}
  This is the declared model definition of force In Newtons.
\end{proof}
\begin{definition}[angular Acceleration In Radians Per Second Squared]
  \label{def:physics:phyx-mini-0778:phyxminiproblems-problemphyxmini0778-angularaccelerationinradianspersecondsquared}
  \lean{PhyXMiniProblems.ProblemPhyXMini0778.angularAccelerationInRadiansPerSecondSquared}
  \uses{def:physics:phyx-mini-0778:phyxminiproblems-problemphyxmini0778-angularaccelerationquantity, def:physics:phyx-mini-0778:phyxminiproblems-problemphyxmini0778-coherentsireadout}
  Radians-per-second-squared readout of angular acceleration
\end{definition}
\begin{proof}
  This is the declared model definition of angular Acceleration In Radians Per Second Squared.
\end{proof}
\begin{definition}[moment Of Inertia In Kilogram Meters Squared]
  \label{def:physics:phyx-mini-0778:phyxminiproblems-problemphyxmini0778-momentofinertiainkilogrammeterssquared}
  \lean{PhyXMiniProblems.ProblemPhyXMini0778.momentOfInertiaInKilogramMetersSquared}
  \uses{def:physics:phyx-mini-0778:phyxminiproblems-problemphyxmini0778-momentofinertiaquantity, def:physics:phyx-mini-0778:phyxminiproblems-problemphyxmini0778-coherentsireadout}
  Kilogram-metre-squared readout of a scalar moment of inertia
\end{definition}
\begin{proof}
  This is the declared model definition of moment Of Inertia In Kilogram Meters Squared.
\end{proof}
\begin{definition}[torque In Newton Meters]
  \label{def:physics:phyx-mini-0778:phyxminiproblems-problemphyxmini0778-torqueinnewtonmeters}
  \lean{PhyXMiniProblems.ProblemPhyXMini0778.torqueInNewtonMeters}
  \uses{def:physics:phyx-mini-0778:phyxminiproblems-problemphyxmini0778-torquequantity, def:physics:phyx-mini-0778:phyxminiproblems-problemphyxmini0778-coherentsireadout}
  Newton-metre readout of an axial torque magnitude
\end{definition}
\begin{proof}
  This is the declared model definition of torque In Newton Meters.
\end{proof}
\begin{definition}[Pictured Body]
  \label{def:physics:phyx-mini-0778:phyxminiproblems-problemphyxmini0778-picturedbody}
  \lean{PhyXMiniProblems.ProblemPhyXMini0778.PicturedBody}
  The three principal bodies visible in the supplied raster
\end{definition}
\begin{proof}
  This is the declared model definition of Pictured Body.
\end{proof}
\begin{definition}[Figure Object]
  \label{def:physics:phyx-mini-0778:phyxminiproblems-problemphyxmini0778-figureobject}
  \lean{PhyXMiniProblems.ProblemPhyXMini0778.FigureObject}
  Individually identifiable physical objects in image 778.png
\end{definition}
\begin{proof}
  This is the declared model definition of Figure Object.
\end{proof}
\begin{definition}[Figure Text Label]
  \label{def:physics:phyx-mini-0778:phyxminiproblems-problemphyxmini0778-figuretextlabel}
  \lean{PhyXMiniProblems.ProblemPhyXMini0778.FigureTextLabel}
  Literal numerical labels printed in the image
\end{definition}
\begin{proof}
  This is the declared model definition of Figure Text Label.
\end{proof}
\begin{definition}[Wire Endpoint]
  \label{def:physics:phyx-mini-0778:phyxminiproblems-problemphyxmini0778-wireendpoint}
  \lean{PhyXMiniProblems.ProblemPhyXMini0778.WireEndpoint}
  The two endpoints of the pictured wire
\end{definition}
\begin{proof}
  This is the declared model definition of Wire Endpoint.
\end{proof}
\begin{definition}[Endpoint Attachment]
  \label{def:physics:phyx-mini-0778:phyxminiproblems-problemphyxmini0778-endpointattachment}
  \lean{PhyXMiniProblems.ProblemPhyXMini0778.EndpointAttachment}
  Bodies to which a wire endpoint can be attached in this setup
\end{definition}
\begin{proof}
  This is the declared model definition of Endpoint Attachment.
\end{proof}
\begin{definition}[Wire Route]
  \label{def:physics:phyx-mini-0778:phyxminiproblems-problemphyxmini0778-wireroute}
  \lean{PhyXMiniProblems.ProblemPhyXMini0778.WireRoute}
  The route of the single wire visible in the primary image
\end{definition}
\begin{proof}
  This is the declared model definition of Wire Route.
\end{proof}
\begin{definition}[Supplied Pulley Figure]
  \label{def:physics:phyx-mini-0778:phyxminiproblems-problemphyxmini0778-suppliedpulleyfigure}
  \lean{PhyXMiniProblems.ProblemPhyXMini0778.SuppliedPulleyFigure}
  \uses{def:physics:phyx-mini-0778:phyxminiproblems-problemphyxmini0778-picturedbody, def:physics:phyx-mini-0778:phyxminiproblems-problemphyxmini0778-figureobject, def:physics:phyx-mini-0778:phyxminiproblems-problemphyxmini0778-figuretextlabel, def:physics:phyx-mini-0778:phyxminiproblems-problemphyxmini0778-wireendpoint, def:physics:phyx-mini-0778:phyxminiproblems-problemphyxmini0778-endpointattachment, def:physics:phyx-mini-0778:phyxminiproblems-problemphyxmini0778-wireroute}
  Typed evidence transcribed from the supplied image. The pulley itself has no mass or diameter label in the raster; those two numerical values come from the problem prose instead
\end{definition}
\begin{proof}
  This is the declared model definition of Supplied Pulley Figure.
\end{proof}
\begin{definition}[Surface Model]
  \label{def:physics:phyx-mini-0778:phyxminiproblems-problemphyxmini0778-surfacemodel}
  \lean{PhyXMiniProblems.ProblemPhyXMini0778.SurfaceModel}
  Model of the horizontal surface described in the prose
\end{definition}
\begin{proof}
  This is the declared model definition of Surface Model.
\end{proof}
\begin{definition}[Wire Mass Model]
  \label{def:physics:phyx-mini-0778:phyxminiproblems-problemphyxmini0778-wiremassmodel}
  \lean{PhyXMiniProblems.ProblemPhyXMini0778.WireMassModel}
  Idealization of the connecting wire stated in the problem
\end{definition}
\begin{proof}
  This is the declared model definition of Wire Mass Model.
\end{proof}
\begin{definition}[Pulley Shape Model]
  \label{def:physics:phyx-mini-0778:phyxminiproblems-problemphyxmini0778-pulleyshapemodel}
  \lean{PhyXMiniProblems.ProblemPhyXMini0778.PulleyShapeModel}
  Geometric mass model of the pulley
\end{definition}
\begin{proof}
  This is the declared model definition of Pulley Shape Model.
\end{proof}
\begin{definition}[Pulley Axle Model]
  \label{def:physics:phyx-mini-0778:phyxminiproblems-problemphyxmini0778-pulleyaxlemodel}
  \lean{PhyXMiniProblems.ProblemPhyXMini0778.PulleyAxleModel}
  Mechanical model of the fixed pulley axle
\end{definition}
\begin{proof}
  This is the declared model definition of Pulley Axle Model.
\end{proof}
\begin{definition}[System State]
  \label{def:physics:phyx-mini-0778:phyxminiproblems-problemphyxmini0778-systemstate}
  \lean{PhyXMiniProblems.ProblemPhyXMini0778.SystemState}
  The instant at which the requested acceleration is evaluated
\end{definition}
\begin{proof}
  This is the declared model definition of System State.
\end{proof}
\begin{definition}[Tabletop Pulley Setup]
  \label{def:physics:phyx-mini-0778:phyxminiproblems-problemphyxmini0778-tabletoppulleysetup}
  \lean{PhyXMiniProblems.ProblemPhyXMini0778.TabletopPulleySetup}
  \uses{def:physics:phyx-mini-0778:phyxminiproblems-problemphyxmini0778-massquantity, def:physics:phyx-mini-0778:phyxminiproblems-problemphyxmini0778-lengthquantity, def:physics:phyx-mini-0778:phyxminiproblems-problemphyxmini0778-accelerationquantity, def:physics:phyx-mini-0778:phyxminiproblems-problemphyxmini0778-forcequantity, def:physics:phyx-mini-0778:phyxminiproblems-problemphyxmini0778-angularaccelerationquantity, def:physics:phyx-mini-0778:phyxminiproblems-problemphyxmini0778-momentofinertiaquantity, def:physics:phyx-mini-0778:phyxminiproblems-problemphyxmini0778-torquequantity, def:physics:phyx-mini-0778:phyxminiproblems-problemphyxmini0778-suppliedpulleyfigure, def:physics:phyx-mini-0778:phyxminiproblems-problemphyxmini0778-surfacemodel, def:physics:phyx-mini-0778:phyxminiproblems-problemphyxmini0778-wiremassmodel, def:physics:phyx-mini-0778:phyxminiproblems-problemphyxmini0778-pulleyshapemodel, def:physics:phyx-mini-0778:phyxminiproblems-problemphyxmini0778-pulleyaxlemodel, def:physics:phyx-mini-0778:phyxminiproblems-problemphyxmini0778-systemstate}
  Independent physical quantities for the two bodies, wire segments, and massive pulley. In particular, boxAcceleration is an independent observable and is not defined from any answer choice or target value
\end{definition}
\begin{proof}
  This is the declared model definition of Tabletop Pulley Setup.
\end{proof}
\begin{definition}[Matches Written Scenario]
  \label{def:physics:phyx-mini-0778:phyxminiproblems-problemphyxmini0778-matcheswrittenscenario}
  \lean{PhyXMiniProblems.ProblemPhyXMini0778.MatchesWrittenScenario}
  \uses{def:physics:phyx-mini-0778:phyxminiproblems-problemphyxmini0778-tabletoppulleysetup}
  Qualitative conditions stated by the problem prose
\end{definition}
\begin{proof}
  This is the declared model definition of Matches Written Scenario.
\end{proof}
\begin{definition}[Matches Primary Figure]
  \label{def:physics:phyx-mini-0778:phyxminiproblems-problemphyxmini0778-matchesprimaryfigure}
  \lean{PhyXMiniProblems.ProblemPhyXMini0778.MatchesPrimaryFigure}
  \uses{def:physics:phyx-mini-0778:phyxminiproblems-problemphyxmini0778-tabletoppulleysetup}
  Geometry, connectivity, and printed labels read directly from image 778.png
\end{definition}
\begin{proof}
  This is the declared model definition of Matches Primary Figure.
\end{proof}
\begin{definition}[Matches Problem Readouts]
  \label{def:physics:phyx-mini-0778:phyxminiproblems-problemphyxmini0778-matchesproblemreadouts}
  \lean{PhyXMiniProblems.ProblemPhyXMini0778.MatchesProblemReadouts}
  \uses{def:physics:phyx-mini-0778:phyxminiproblems-problemphyxmini0778-massinkilograms, def:physics:phyx-mini-0778:phyxminiproblems-problemphyxmini0778-lengthinmeters, def:physics:phyx-mini-0778:phyxminiproblems-problemphyxmini0778-tabletoppulleysetup}
  Physical numerical data from the two image labels and the problem prose. 0.500 m is represented exactly as 1 / 2 m. The radius is not assigned here; it is related to the stated diameter by a general disk-geometry law
\end{definition}
\begin{proof}
  This is the declared model definition of Matches Problem Readouts.
\end{proof}
\begin{definition}[Uses Standard Near Earth Gravity]
  \label{def:physics:phyx-mini-0778:phyxminiproblems-problemphyxmini0778-usesstandardnearearthgravity}
  \lean{PhyXMiniProblems.ProblemPhyXMini0778.UsesStandardNearEarthGravity}
  \uses{def:physics:phyx-mini-0778:phyxminiproblems-problemphyxmini0778-accelerationinmeterspersecondsquared, def:physics:phyx-mini-0778:phyxminiproblems-problemphyxmini0778-tabletoppulleysetup}
  The gravitational calibration implicit in the recorded textbook answer, separated from the values explicitly printed in the problem
\end{definition}
\begin{proof}
  This is the declared model definition of Uses Standard Near Earth Gravity.
\end{proof}
\begin{definition}[Has Physical Parameters]
  \label{def:physics:phyx-mini-0778:phyxminiproblems-problemphyxmini0778-hasphysicalparameters}
  \lean{PhyXMiniProblems.ProblemPhyXMini0778.HasPhysicalParameters}
  \uses{def:physics:phyx-mini-0778:phyxminiproblems-problemphyxmini0778-massinkilograms, def:physics:phyx-mini-0778:phyxminiproblems-problemphyxmini0778-lengthinmeters, def:physics:phyx-mini-0778:phyxminiproblems-problemphyxmini0778-accelerationinmeterspersecondsquared, def:physics:phyx-mini-0778:phyxminiproblems-problemphyxmini0778-forceinnewtons, def:physics:phyx-mini-0778:phyxminiproblems-problemphyxmini0778-momentofinertiainkilogrammeterssquared, def:physics:phyx-mini-0778:phyxminiproblems-problemphyxmini0778-tabletoppulleysetup}
  Positivity and nondegeneracy conditions for the massive-pulley model
\end{definition}
\begin{proof}
  This is the declared model definition of Has Physical Parameters.
\end{proof}
\begin{definition}[Satisfies Massive Pulley Dynamics]
  \label{def:physics:phyx-mini-0778:phyxminiproblems-problemphyxmini0778-satisfiesmassivepulleydynamics}
  \lean{PhyXMiniProblems.ProblemPhyXMini0778.SatisfiesMassivePulleyDynamics}
  \uses{def:physics:phyx-mini-0778:phyxminiproblems-problemphyxmini0778-massinkilograms, def:physics:phyx-mini-0778:phyxminiproblems-problemphyxmini0778-lengthinmeters, def:physics:phyx-mini-0778:phyxminiproblems-problemphyxmini0778-accelerationinmeterspersecondsquared, def:physics:phyx-mini-0778:phyxminiproblems-problemphyxmini0778-forceinnewtons, def:physics:phyx-mini-0778:phyxminiproblems-problemphyxmini0778-angularaccelerationinradianspersecondsquared, def:physics:phyx-mini-0778:phyxminiproblems-problemphyxmini0778-momentofinertiainkilogrammeterssquared, def:physics:phyx-mini-0778:phyxminiproblems-problemphyxmini0778-torqueinnewtonmeters, def:physics:phyx-mini-0778:phyxminiproblems-problemphyxmini0778-tabletoppulleysetup}
  The ideal-wire and massive-disk-pulley equations at the instant after release. All signs are fixed by taking motion of the box toward the pulley, motion of the hanging weight downward, and pulley rotation in the corresponding sense as positive. The two wire tensions remain independent because a massive pulley requires a nonzero tension difference. None of these general laws mentions 49 / 18, 2.72, or an answer label
\end{definition}
\begin{proof}
  This is the declared model definition of Satisfies Massive Pulley Dynamics.
\end{proof}
\begin{lemma}[box Acceleration exact]
  \label{lem:physics:phyx-mini-0778:phyxminiproblems-problemphyxmini0778-boxacceleration-exact}
  \lean{PhyXMiniProblems.ProblemPhyXMini0778.boxAcceleration_exact}
  \uses{def:physics:phyx-mini-0778:phyxminiproblems-problemphyxmini0778-accelerationinmeterspersecondsquared, def:physics:phyx-mini-0778:phyxminiproblems-problemphyxmini0778-tabletoppulleysetup, def:physics:phyx-mini-0778:phyxminiproblems-problemphyxmini0778-matcheswrittenscenario, def:physics:phyx-mini-0778:phyxminiproblems-problemphyxmini0778-matchesprimaryfigure, def:physics:phyx-mini-0778:phyxminiproblems-problemphyxmini0778-matchesproblemreadouts, def:physics:phyx-mini-0778:phyxminiproblems-problemphyxmini0778-usesstandardnearearthgravity, def:physics:phyx-mini-0778:phyxminiproblems-problemphyxmini0778-hasphysicalparameters, def:physics:phyx-mini-0778:phyxminiproblems-problemphyxmini0778-satisfiesmassivepulleydynamics}
  The governing equations yield a = m\_h g / (m\_box + m\_h + I / r²). For a uniform disk, I / r² = m\_pulley / 2, so the supplied data give the exact acceleration 49 / 18 m/s²
\end{lemma}
\begin{proof}
  The conclusion follows from the governing relations and figure readouts named in the statement of box Acceleration exact.
\end{proof}
\begin{definition}[Answer Choice]
  \label{def:physics:phyx-mini-0778:phyxminiproblems-problemphyxmini0778-answerchoice}
  \lean{PhyXMiniProblems.ProblemPhyXMini0778.AnswerChoice}
  Labels of the four answers printed in the problem
\end{definition}
\begin{proof}
  This is the declared model definition of Answer Choice.
\end{proof}
\begin{definition}[acceleration In Meters Per Second Squared]
  \label{def:physics:phyx-mini-0778:phyxminiproblems-problemphyxmini0778-answerchoice-accelerationinmeterspersecondsquared}
  \lean{PhyXMiniProblems.ProblemPhyXMini0778.AnswerChoice.accelerationInMetersPerSecondSquared}
  \uses{def:physics:phyx-mini-0778:phyxminiproblems-problemphyxmini0778-accelerationinmeterspersecondsquared, def:physics:phyx-mini-0778:phyxminiproblems-problemphyxmini0778-answerchoice}
  Acceleration printed beside each answer label, in m/s²
\end{definition}
\begin{proof}
  This is the declared model definition of acceleration In Meters Per Second Squared.
\end{proof}
\begin{definition}[recorded Answer Choice]
  \label{def:physics:phyx-mini-0778:phyxminiproblems-problemphyxmini0778-recordedanswerchoice}
  \lean{PhyXMiniProblems.ProblemPhyXMini0778.recordedAnswerChoice}
  \uses{def:physics:phyx-mini-0778:phyxminiproblems-problemphyxmini0778-answerchoice}
  The answer label recorded in the supplied dataset metadata
\end{definition}
\begin{proof}
  This is the declared model definition of recorded Answer Choice.
\end{proof}
\begin{definition}[Matches Displayed Answer]
  \label{def:physics:phyx-mini-0778:phyxminiproblems-problemphyxmini0778-matchesdisplayedanswer}
  \lean{PhyXMiniProblems.ProblemPhyXMini0778.MatchesDisplayedAnswer}
  \uses{def:physics:phyx-mini-0778:phyxminiproblems-problemphyxmini0778-accelerationinmeterspersecondsquared, def:physics:phyx-mini-0778:phyxminiproblems-problemphyxmini0778-tabletoppulleysetup, def:physics:phyx-mini-0778:phyxminiproblems-problemphyxmini0778-answerchoice}
  Agreement with an acceleration printed to two decimal places. The tolerance 0.005 m/s² = 1/200 m/s² is half a unit in the last displayed digit, so this does not falsely identify the unrounded value 49 / 18 with the decimal 2.72
\end{definition}
\begin{proof}
  This is the declared model definition of Matches Displayed Answer.
\end{proof}
% --- Archon declaration topology end ---
% --- Archon physics formalization source end ---
